\subsection{Nomenclatura de Métodos de Prueba}

\subsubsection{Regla: patrón método-estado-comportamiento esperado}

\textbf{Regla:} Nombre en tres partes, en \texttt{PascalCase} separadas por guion bajo: método bajo prueba, estado evaluado, comportamiento esperado (Osherove, 2013). Esta es una excepción explícita a la regla general de nombres de método sin separadores. Los ejemplos omiten el atributo de prueba porque la regla de nombrado aplica al margen de él.

\textbf{Código correcto:}
\begin{lstlisting}[style=csharpstyle]
public void Roll_FaceValueSix_GrantsExtraTurn()
{
}
\end{lstlisting}

\textbf{Código incorrecto:}
\begin{lstlisting}[style=csharpstyle]
public void TestRollExtraTurn()
{
}
\end{lstlisting}
